\documentclass[conference]{IEEEtran}
\IEEEoverridecommandlockouts
\usepackage{cite}
\usepackage{amsmath,amssymb,amsfonts}
\usepackage{algorithmic}
\usepackage{graphicx}
\usepackage{textcomp}
\usepackage{xcolor}
\usepackage{booktabs}
\usepackage{multirow}
\usepackage{algorithm}
\usepackage{array}
\usepackage{url}
\usepackage{balance}
\usepackage{subcaption}

\begin{document}

\title{Mining Experiment-Similarity Graphs from Code\\for Cross-Paper Hypothesis Generation}

\author{
\IEEEauthorblockN{Divyansh Bhatia}
\IEEEauthorblockA{eBay Inc.\\dnbhatia@ebay.com}
\and
\IEEEauthorblockN{Dhiyanesh Sidhaiyan}
\IEEEauthorblockA{ACV Auctions Inc.\\dsidhaiyan@acvauctions.com}
}

\maketitle

\begin{abstract}
Machine learning researchers collectively explore vast experimental spaces---ablation studies, benchmark evaluations, hyperparameter sweeps---yet cross-paper synthesis of these findings remains manual and ad hoc.
We introduce a system that mines structured experimental features from 12{,}303 code repositories spanning 28 venues, constructs an experiment-similarity graph over 2{,}530 papers, and applies community detection with coherence-based filtering to identify cross-venue research communities.
Scoped LLM reasoning over each community generates cross-paper hypotheses grounded in members' actual experimental evidence.
At similarity threshold $\tau$=6.0, our method produces 50 hypotheses scoring 3.79/5.0 mean quality across specificity, groundedness, actionability, and cross-paper novelty---outperforming standalone LLM generation (+0.66), RAG retrieval (+0.32), and abstract-keyword graphs (+1.00), while achieving 93\% of AI~Scientist-style refinement quality at 10$\times$ the hypothesis volume.
A controlled baseline analysis with two judge models reveals that community scoping accounts for the primary gain (+0.22 under a stronger judge); structured prompting adds specificity (+1.40) but not mean quality, while also surfacing a methodological finding: LLM-as-judge rankings can invert depending on the judge's ability to detect hallucinated evidence.
A threshold ablation over $\tau \in [2.0, 6.0]$ shows that tight communities sharing concrete benchmarks yield highest quality (3.75), and an LLM ablation across 7~models (5~runs each, with statistical significance) demonstrates that the pipeline is model-agnostic: even Llama~3.3~70B achieves 139.4 volume$\times$quality, surpassing all non-graph baselines.
Per-paper ablation methods score 1.0/5.0 on cross-paper novelty, confirming that graph structure is essential.
We validate 8 system-generated hypotheses as actual experiments across 4~domains: 4 validate, 3 partially validate with overestimated effect sizes, and 1 fails due to architectural constraints invisible to the system.
\end{abstract}

\begin{IEEEkeywords}
experiment mining, hypothesis generation, graph mining, community detection, scientific discovery
\end{IEEEkeywords}

%% ====================================================================
\section{Introduction}
%% ====================================================================

Machine learning research produces thousands of papers annually, each contributing experimental evidence---benchmark results, ablation studies, hyperparameter analyses---to a collective but fragmented knowledge base.
A CVPR paper showing that increasing image resolution from 256 to 320 pixels improves anomaly localization by 3--5\% and a NeurIPS paper demonstrating that batch size optimization within $\{12, 32, 64\}$ affects convergence on the same benchmark are \emph{experimentally related}: combining their findings suggests a specific untested experiment.
Yet no existing system surfaces such cross-paper experimental connections at scale.

Current approaches to automated hypothesis generation operate either on paper text~\cite{si2024llmideas, radensky2025scideator, lu2024aiscientist} or on individual papers in isolation~\cite{abramovich2025ablationbench, zhao2025abgen}.
Text-based methods capture topical similarity but miss the structured experimental data---specific parameter values, benchmark scores, ablation configurations---that resides in code repositories.
Per-paper methods produce actionable within-paper suggestions but cannot synthesize findings across the literature.
Knowledge-graph approaches~\cite{xiong2024kgcoi, ghafarollahi2024sciagents} build graphs from text-extracted entities but lack the experimental grounding needed for specific, testable hypotheses.

We address this gap with a four-stage pipeline (Fig.~\ref{fig:pipeline}).
First, we mine structured experimental features from code repositories via static analysis, processing 12{,}303 repositories across 28 venues covering ML, CV, NLP, and data mining.
Second, we construct an \emph{experiment-similarity graph} where edges encode shared experimental configurations, using inverted-index batching for scalability (91\% reduction in pairwise comparisons).
Third, we apply Louvain community detection with coherence-based clique refinement to identify cross-venue research communities.
Fourth, we scope LLM reasoning to each community to generate hypotheses grounded in members' actual metric values.

We make the following contributions:
\begin{itemize}
\item We formalize cross-paper hypothesis generation as community detection in experiment-similarity graphs and propose a scalable four-stage pipeline processing 12{,}303 repositories from 28 venues (Section~III).
\item We evaluate against 5~baselines on 4~dimensions, showing +0.66 over standalone LLM and +1.00 over abstract-keyword graphs, achieving 93\% of AI~Scientist quality at 10$\times$ the volume (Section~V-A).
\item We conduct a threshold ablation ($\tau \in [2.0, 6.0]$) revealing that tight communities with shared benchmarks yield highest quality (3.75), and a rank degradation analysis validating the scoring function (Section~V-B).
\item We perform the first cross-model ablation for hypothesis generation: 7~LLMs $\times$ 5~runs with paired $t$-tests ($p < 0.001$) and Kruskal-Wallis tests, showing the pipeline is model-agnostic (Section~V-C).
\item We validate 8 system-generated hypotheses as actual experiments across CV, graph mining, NLP, and PEFT---achieving 87.5\% directional accuracy---and verify novelty through literature search (80\% novel, 0\% known) and human evaluation with 6 ML researchers (Sections~IV-C, V-K).
\end{itemize}


%% ====================================================================
%% PIPELINE FIGURE (two-column)
%% ====================================================================
\begin{figure*}[t]
\centerline{\includegraphics[width=\textwidth]{fig1_pipeline_clean.pdf}}
\caption{System pipeline. The system extracts structured experimental features from 12{,}303 repositories across 28 venues. The experiment-similarity graph connects papers sharing benchmarks, methods, or configurations. Community detection identifies cross-venue groups sharing concrete benchmarks. LLM reasoning within each community generates cross-paper hypotheses grounded in specific experimental evidence. Italic labels trace the narrative: \emph{What experiments did they run?} $\rightarrow$ \emph{Which papers share setups?} $\rightarrow$ \emph{Same discussion group?} $\rightarrow$ \emph{What has nobody tried?}}
\label{fig:pipeline}
\end{figure*}


%% ====================================================================
\section{Related Work}
%% ====================================================================

\subsection{LLM-Driven Scientific Idea Generation}

Si et al.~\cite{si2024llmideas} studied LLM-generated research ideas with 100+ NLP researchers, finding LLM ideas were rated more novel but less feasible than human ideas.
The AI Scientist~\cite{lu2024aiscientist} automates the full research lifecycle---ideation, experimentation, paper writing---and achieved workshop-level acceptance~\cite{yamada2025aiscientistv2}.
Scideator~\cite{radensky2025scideator} grounds generation in ``facet recombination,'' decomposing papers into purpose, mechanism, and evaluation facets.
HypoAgents~\cite{hypoagents2025} uses multi-agent Bayesian-entropy reasoning for hypothesis generation.
ResearchAgent~\cite{baek2024researchagent} iteratively refines ideas using academic graph knowledge, while SciMON~\cite{wang2024scimon} optimises generated ideas for novelty against the existing literature.
Recent benchmarks~\cite{researchbench2025} evaluate LLM idea quality against ground-truth, but none mine structured experimental features from code.

\subsection{Experiment and Ablation Mining}

AblationBench~\cite{abramovich2025ablationbench} evaluates LLMs on planning ablation studies from method sections.
AbGen~\cite{zhao2025abgen} generates ablation designs from method descriptions.
Both produce actionable per-paper suggestions but cannot synthesize across papers.
MLCopilot~\cite{zhang2024mlcopilot} leverages historical ML experiments to suggest configurations for new tasks, but operates on curated benchmarks rather than mined repositories.
Our extraction system mines code repositories---YAML configurations, argparse defaults, shell scripts, README result tables---and extends per-repository insights into a cross-paper graph.
Experiment tracking platforms (OpenML, MLflow, Weights~\&~Biases) offer complementary data; integrating these is a natural extension.

\subsection{Knowledge Graphs for Scientific Discovery}

Swanson's literature-based discovery~\cite{swanson1986} demonstrated that connecting disjoint literatures yields novel hypotheses---the ``ABC model'' where A relates to B and B to C, but the A--C connection is undiscovered.
Modern systems build graphs from text-extracted entities: KG-CoI~\cite{xiong2024kgcoi} uses chain-of-ideas reasoning, SciAgents~\cite{ghafarollahi2024sciagents} employs multi-agent graph reasoning, and Tong et al.~\cite{tong2024causal} use causal graphs.
Our experiment-similarity graph differs: edges encode \emph{shared experimental configurations} from code, not textual co-occurrence---a relationship invisible to text-based methods but immediately actionable for experiment design.


%% ====================================================================
\section{Methodology}
%% ====================================================================

\subsection{Problem Formulation}

\begin{definition}[Experimental Feature Set]
For paper $p$ with repository $r$, the feature set $F_p = (B_p, M_p, P_p, A_p, \mu_p)$ consists of benchmarks $B_p$, method keywords $M_p$, hyperparameters $P_p$, ablated parameters $A_p$, and metrics $\mu_p$, extracted from $r$ via static analysis.
The first four components drive graph construction (Equation~1); metrics $\mu_p$ are reserved for hypothesis grounding in Stage~4.
\end{definition}

\begin{definition}[Experiment-Similarity Graph]
An undirected weighted graph $G = (V, E, w)$ where $V$ is the set of papers and edge $(v_i, v_j) \in E$ with weight $w(v_i, v_j) \geq \tau$ encodes shared experimental configurations.
\end{definition}

\noindent\textbf{Problem.} Given paper--repository pairs $\{(p_i, r_i)\}_{i=1}^{n}$ and threshold $\tau$, construct $G$, detect communities $\{C_1, \ldots, C_k\}$, and for each $C_j$ generate cross-paper hypotheses grounded in members' experimental evidence.


\subsection{Stage 1: Code-Level Feature Extraction}

Given a repository $r$, we extract features through static analysis (no execution required).
The extraction pipeline processes four file types:

\textbf{Configuration files.}
We parse YAML, JSON, and TOML files containing $\geq$2 ML-relevant keywords (e.g., \texttt{learning\_rate}, \texttt{batch\_size}, \texttt{backbone}, \texttt{dataset}).
Across 12{,}303 repos, this yields 117{,}428 config files---the dominant source at 77\% of all extracted files.

\textbf{Shell scripts.}
Launch scripts (e.g., \texttt{train.sh}, \texttt{run\_experiments.sh}) are parsed for flag--value patterns (\texttt{-{}-lr 0.001}, \texttt{-{}-backbone resnet50}).
These contribute 35{,}537 files (23\%), often encoding ablation sweeps as parallel commands with varying arguments.

\textbf{Python scripts.}
\texttt{argparse} definitions are parsed to extract parameter names, types, and defaults.
This captures experimental parameters not specified in config files.

\textbf{Result tables.}
Markdown tables in README files with $\geq$30\% numerical cells are parsed to extract model--metric pairs (e.g., ``ViT-B: 82.1 mAP'').

\textbf{Canonicalization.}
Parameters are canonicalized via a mapping of 45~aliases to 17~groups (e.g., \texttt{lr}, \texttt{learning\_rate}, \texttt{optim.lr} $\rightarrow$ \texttt{learning\_rate}).
Benchmarks are detected via keyword matching against 27~known benchmarks.
From these raw extractions, we construct per-repository \emph{experiment graphs} identifying seven insight types (Table~\ref{tab:extraction}): parameter ranges, clean ablations (single parameter varied), experiment families (shared structure with varying configs), result comparisons, cross-dataset ablations, configuration--result links, and dataset comparisons.


\subsection{Stage 2: Graph Construction}

For papers $v_i$, $v_j$ with features $F_i$, $F_j$, the edge weight is:
\begin{equation}
w(v_i, v_j) = 3|B_i \!\cap\! B_j| + 2|M_i \!\cap\! M_j| + \alpha|P_i \!\cap\! P_j| + 2.5|A_i \!\cap\! A_j|
\end{equation}
where $\alpha = 1.0$ for domain-specific parameters and $0.3$ for common ones (learning rate, batch size, epochs, seed) to avoid trivially connecting all papers.
The weights reflect experimental relevance: shared benchmarks ($\times$3.0) indicate papers that can directly compare results; shared ablation parameters ($\times$2.5) indicate papers exploring the same experimental axes; shared methods ($\times$2.0) indicate methodological overlap.
Common parameters receive reduced weight ($\times$0.3) because nearly all ML papers share learning rate and batch size.
An edge is added when $w \geq \tau$.

\textbf{Inverted-index batching.}
Na\"ive pairwise computation over 2{,}530 papers requires 3.2M comparisons.
We build inverted indices mapping each feature to its containing papers (benchmark $\rightarrow$ papers, method $\rightarrow$ papers, parameter $\rightarrow$ papers).
Similarity is computed only for pairs sharing $\geq$1 feature, reducing candidates to 280{,}728 (91.2\% reduction).
At $\tau$=6.0 this produces 16{,}027 edges; at $\tau$=2.0, 280{,}408 edges.


\subsection{Stage 3: Community Detection}

We apply Louvain partitioning~\cite{blondel2008louvain} at resolution 1.2 to maximize modularity~\cite{fortunato2010community}.
Raw Louvain communities can be large and heterogeneous; we refine communities with $\geq$8 members by extracting maximal cliques of size $\geq$4, producing tighter subgroups.
Each community $C$ is scored:
\begin{equation}
\text{score}(C) = \text{coh}(C) \cdot (1 + 0.3|V_C|) \cdot \mathbb{1}_{\text{xdom}} \cdot 1.5
\end{equation}
where $\text{coh}(C) = 3|\cap_i B_i| + 2|\cap_i M_i| + 1.5|\cap_i A_i|$ measures shared experimental features across all members, $|V_C|$ is venue diversity (rewarding communities drawing from multiple conferences), and $\mathbb{1}_{\text{xdom}}$ is a bonus for communities spanning multiple research domains, defined by venue: CV (CVPR, ECCV), NLP (ACL, EMNLP, NAACL), ML (ICLR, NeurIPS, ICML, UAI), DM (KDD, ICDM, AAAI).
Communities are deduplicated by removing those with $>$50\% member overlap with higher-ranked ones.
The scoring function prioritizes communities where (a)~all members share a concrete benchmark, (b)~papers come from diverse venues, and (c)~multiple domains are represented.
At $\tau$=6.0, this produces 10 communities with 100\% shared benchmarks (Section~V-B).
Community detection is highly stable: across 10 random seeds, Louvain yields NMI = 0.998$\pm$0.001 and ARI = 0.959$\pm$0.022, confirming that the community structure is determined by the graph topology, not initialization.
Communities are also stable \emph{across} $\tau$ values: NMI between $\tau$=4.0 and $\tau$=6.0 is 0.929, with 43\% of $\tau$=6.0 communities being strict subsets of $\tau$=4.0 communities---increasing $\tau$ refines rather than restructures the community landscape.


\subsection{Stage 4: Hypothesis Generation}

For each community, we construct a \emph{metric-grounded prompt} (Fig.~\ref{fig:prompt}) containing three components.
First, \textbf{member evidence}: for each paper in the community, we list its benchmarks, methods, and up to 6 key insights.
Insights with quantitative evidence (config--result links, clean ablations) are formatted with $\bigstar$ markers showing actual metric values---e.g., ``$\bigstar$ ResNet-50 = 78.3 mAP, ViT-B = 82.1 mAP.''
This ensures the LLM has concrete numbers to cite, not just textual descriptions.
Second, \textbf{community profile}: shared benchmarks, all benchmarks tested (with coverage counts), ablated parameters, and key metrics.
Third, \textbf{grounding rules} requiring the LLM to (a)~cite specific metric values from $\geq$2 papers, (b)~predict quantitative expected outcomes, and (c)~omit any hypothesis lacking numerical support.

The LLM generates 5 hypotheses per community in structured JSON format, each containing: hypothesis statement, motivation with cited numbers, expected outcome, confidence level, cited papers, and gap type (unexplored combination, cross-venue transfer, untested transfer, or missing ablation).
Algorithm~\ref{alg:pipeline} summarizes the complete pipeline.


\begin{figure}[t]
\centering
\footnotesize
\fbox{\parbox{0.93\columnwidth}{%
\textbf{System prompt} (abbreviated):\\[2pt]
\textit{You are a research advisor analyzing a cluster of \textnormal{N} ML papers from \textnormal{V} venues. These papers share \textnormal{S} benchmarks.}\\[4pt]
\textbf{Member evidence:}\\
\texttt{[Paper: "Efficient ViT Pruning..."]}\\
\texttt{\ \ Benchmarks: [CIFAR-10, CIFAR-100]}\\
\texttt{\ \ Methods: [vit, pruning]}\\
\texttt{\ \ }$\bigstar$ \texttt{pruning\_ratio=0.5: acc=94.2}\\
\texttt{\ \ }$\bigstar$ \texttt{pruning\_ratio=0.7: acc=93.9}\\[2pt]
\texttt{[Paper: "Curriculum Learning for..."]}\\
\texttt{\ \ Benchmarks: [CIFAR-100, ImageNet]}\\
\texttt{\ \ Methods: [resnet, curriculum]}\\
\texttt{\ \ }$\bigstar$ \texttt{curriculum: +3.2\% accuracy}\\
\texttt{\ \ }$\bigstar$ \texttt{no\_curriculum: 76.1\% accuracy}\\[4pt]
\textbf{Grounding rules:}\\
\textit{-- Every hypothesis MUST cite} $\geq$\textit{1 specific}\\
\textit{\ \ \ metric value from each of 2+ papers.}\\
\textit{-- Use actual numbers from} $\bigstar$ \textit{lines.}\\
\textit{-- Predict QUANTITATIVE outcomes.}\\
\textit{-- If no numbers support it, do NOT include it.}\\[4pt]
\textbf{Output:} JSON array of 5 hypotheses.
}}
\caption{Abbreviated prompt template for community-grounded hypothesis generation.
Member evidence includes actual metric values (marked $\bigstar$) extracted from code repositories.
Grounding rules enforce citation of specific numbers.}
\label{fig:prompt}
\end{figure}


\begin{algorithm}[t]
\caption{Cross-Paper Hypothesis Generation}
\label{alg:pipeline}
\begin{algorithmic}[1]
\REQUIRE Corpus $\mathcal{D} = \{(p_i, r_i)\}$, threshold $\tau$
\ENSURE Hypotheses $\mathcal{H}$
\FOR{$(p_i, r_i) \in \mathcal{D}$}
    \STATE $F_i \leftarrow \textsc{ExtractFeatures}(r_i)$
\ENDFOR
\STATE $G \leftarrow \textsc{BuildGraph}(\{F_i\}, \tau)$ \hfill\COMMENT{inverted index}
\STATE $\{C_j\} \leftarrow \textsc{Communities}(G)$ \hfill\COMMENT{Louvain + refine}
\STATE Rank by $\text{score}(C_j)$; deduplicate
\FOR{$C_j \in \text{top-}K$ communities}
    \STATE $\mathcal{H} \leftarrow \mathcal{H} \cup \textsc{LLM}(\textsc{FormatEvidence}(C_j))$
\ENDFOR
\RETURN $\mathcal{H}$
\end{algorithmic}
\end{algorithm}


%% ====================================================================
\section{Experimental Setup}
%% ====================================================================

\subsection{Dataset}

We collected 28{,}271 papers from 28 venues via OpenReview API (ICLR 2024--2025, NeurIPS 2023--2024, ICML 2023--2024, UAI 2024) and PapersWithCode data dumps (CVPR, ECCV, ACL, EMNLP, NAACL, AAAI, KDD, ICDM), yielding 12{,}303 with linked code repositories.
By domain: ML venues contribute 5{,}904 repos (48\%), CV 4{,}453 (36\%), NLP 1{,}600 (13\%), and DM/AI 106 (1\%).
Feature extraction yielded 2{,}530 papers with sufficient features for graph construction (Table~\ref{tab:extraction}).
The top benchmarks are COCO (201 papers), ImageNet (172), VOC (128), CIFAR-10 (97); top methods are ViT (231), pretraining (203), CLIP (200), ResNet (190).


\subsection{Baselines}

We compare against five approaches, all using the same frontier LLM and receiving the same extracted insights:

\textbf{B1:~Standalone LLM}~\cite{si2024llmideas}---papers' insights fed in batches without structural grouping.
\textbf{B2:~RAG+LLM}~\cite{radensky2025scideator}---for seed papers, most similar are retrieved by keyword overlap.
\textbf{B3:~Per-paper ablation}~\cite{abramovich2025ablationbench}---each paper analyzed individually.
\textbf{B4:~Abstract-keyword graph}~\cite{xiong2024kgcoi}---same community detection and prompt, edges from title keywords.
\textbf{B5:~AI Scientist-style}~\cite{lu2024aiscientist}---two-step ideation with self-refinement.
Each baseline's inherent design determines how many papers it processes (B1: 100 in batches, B3: 20 individually, B5: 60 in context), isolating the effect of graph-based community scoping.


\subsection{Evaluation Protocol}

Hypotheses from all methods are shuffled into a blind pool and rated by an LLM judge (Model~A\footnote{Model~A is Claude Sonnet~4 (Anthropic); Model~B is Claude Opus~4.5 (Anthropic). Model~B has ${\sim}3\times$ more parameters and scores higher on MMLU/HumanEval benchmarks. A different-provider judge (GPT-5.2) is also tested; see Section~V-C.}) on four dimensions (1--5):
\textbf{Specificity}~(Spec): exact parameters, values, benchmarks named.
\textbf{Groundedness}~(Grnd): backed by specific metric values from cited papers.
\textbf{Actionability}~(Actn): implementable as a concrete experiment.
\textbf{Cross-paper novelty}~(XPap): requires knowledge from $\geq$2 papers; single-paper methods should score ${\sim}$1.
We define volume$\times$quality (V${\times}$Q) as the product of hypothesis count and mean quality, measuring total productive output; we note this assumes linear value in count, which overestimates marginal contributions at high volumes.

\textbf{Judge reliability.}
To validate LLM-as-judge, we had two independent judges (Model~A and Model~B, from different providers) rate the same 49 hypotheses on all four dimensions.
Quadratic-weighted Cohen's $\kappa = 0.52$ (moderate agreement), with 83\% of ratings within 1~point.
Kendall's $\tau = 0.63$ ($p < 0.001$) and Spearman's $\rho = 0.72$ confirm strong rank correlation.
Cross-paper novelty shows the highest agreement ($\kappa = 0.71$, substantial).
Model~B is systematically stricter (mean 2.69 vs.\ Model~A's 3.52), particularly on groundedness (1.61 vs.\ 2.84), but the relative ordering of hypotheses is preserved ($\tau = 0.52$, $p < 0.001$).
Since our evaluations compare methods by \emph{relative} quality, not absolute scores, this rank preservation is the critical property.
Furthermore, if judge bias favored its own model's generation style, the same model used as generator should rank highest in the LLM ablation---but it ranks 5th of 7 (Table~\ref{tab:llm}), behind four models from different providers.
Under the stricter Model~B judge, our method's 45 hypotheses score 3.04 mean---still above the abstract-keyword baseline (B4, 2.79), confirming that our results hold under independent evaluation.

\textbf{Human evaluation.}
Six ML researchers (4--8 years experience) rated 10 randomly sampled hypotheses on scientific plausibility (3.0/5), practical value (2.7/5), and feasibility (3.5/5), with moderate inter-rater agreement (Kendall's $\tau = 0.58$).
Four of six said they would use the system; dissenting evaluators flagged the lack of mechanistic reasoning.
Hypotheses identifying benchmarking gaps (H9: 3.7/5) and augmentation transfers (H3: 3.8/5) received the highest ratings, while superficial cross-paper parameter matching (H5: 2.2/5, H7: 2.4/5) scored lowest.
The two most critical evaluators (overall scores 2.0/5) specifically noted: ``the system often notices that Paper~A tuned parameter~X and Paper~B also has parameter~X, then assumes the transfer is meaningful---that is often not enough'' and ``no mechanistic reasoning for why a parameter change would produce the predicted effect.''
Evaluators consistently flagged overestimated effect sizes and miscalibrated confidence as primary failure modes---consistent with our calibration analysis showing 31\% of predictions exceed 15\%.
We confirmed one such failure by running a top-rated hypothesis (resolution ablation for RealNet~\cite{zhang2024realnet}): the model's UNet skip connections are hardcoded for 256$\times$256, making the experiment architecturally impossible---a constraint invisible to config-level extraction.
Three of ten hypotheses were rated strong ($\geq$3.5), three okay (3.0--3.5), and four weak ($<$3.0), suggesting the system functions best as an idea generator requiring expert triage.


%% ====================================================================
\section{Results and Analysis}\label{sec:results}
%% ====================================================================

\subsection{Baseline Comparison}

Table~\ref{tab:main_results} presents the main comparison at $\tau$=6.0.
Fig.~\ref{fig:baselines} visualizes both per-dimension quality and volume-adjusted productivity.

\begin{table}[t]
\caption{Baseline comparison ($\tau$=6.0). Bold = best; underline = second. V$\times$Q = volume $\times$ mean quality. See Section~VI for controlled decomposition under a stronger judge.}
\label{tab:main_results}
\centering
\footnotesize
\setlength{\tabcolsep}{3pt}
\begin{tabular}{lcccccrr}
\toprule
\textbf{Method} & \textbf{Spec} & \textbf{Grnd} & \textbf{Actn} & \textbf{XPap} & \textbf{Mean} & $n$ & \textbf{V$\times$Q} \\
\midrule
\textbf{Ours} & \underline{4.07} & \underline{3.40} & \underline{4.07} & 3.60 & \underline{3.79} & \textbf{50} & \textbf{189.5} \\
B5: AI Sci. & \textbf{4.60} & \textbf{3.80} & 3.60 & \textbf{4.20} & \textbf{4.05} & 5 & 20.2 \\
B2: RAG & 3.60 & 3.27 & 3.33 & \underline{3.67} & 3.47 & 30 & 104.1 \\
B1: Standalone & 3.67 & 2.73 & 3.21 & 2.93 & 3.13 & 40 & 125.2 \\
B3: Ablation & \underline{4.07} & 2.80 & \textbf{4.67} & 1.00 & 3.13 & 40 & 125.2 \\
B4: Abstr.~graph & 2.67 & 2.73 & 2.47 & 3.27 & 2.79 & 30 & 83.7 \\
\bottomrule
\end{tabular}
\end{table}


\begin{figure}[t]
\centering
\includegraphics[width=\columnwidth]{fig2_baselines.pdf}
\caption{(a)~Hypothesis quality by dimension. Our method achieves balanced strength across all four metrics. B3 wins on actionability but scores 1.0 on cross-paper novelty. (b)~Volume$\times$quality: our method achieves 93\% of B5's per-hypothesis quality at 10$\times$ the volume.}
\label{fig:baselines}
\end{figure}


\textbf{Finding 1: Code features outperform text features.}
Our approach outperforms B4 (abstract-keyword graph) by +1.00 mean, with the largest gaps on specificity (+1.40) and actionability (+1.60).
To test a stronger text baseline, we replaced code features with Sentence-BERT~\cite{reimers2019sbert} title embeddings (cosine threshold 0.50 for matched edge count).
Embedding communities achieve 2.00 coherence with 0/9 shared benchmarks vs.\ code-feature communities at 5.06 coherence with 8/8 shared benchmarks---a 2.5$\times$ coherence gap.
Hypothesis quality is also higher for code features (3.92 vs.\ 3.71).
Text similarity groups topically related papers; code features group \emph{experimentally} related papers that share concrete benchmarks and configurations.

\textbf{Finding 2: Graph structure improves over flat LLM.}
We outperform B1 (standalone LLM) on all four dimensions: specificity (+0.40), groundedness (+0.67), actionability (+0.86), cross-paper novelty (+0.67).

\textbf{Finding 3: Cross-paper novelty proves the graph is essential.}
B3 (per-paper ablation) scores 1.00 on cross-paper novelty---the minimum possible.
It excels at actionability (4.67) because single-paper ablations are immediately implementable, but it \emph{cannot} synthesize across papers.
Our 3.60 vs.\ B3's 1.00 confirms that graph structure enables a qualitatively different class of hypotheses.

\textbf{Finding 4: Best volume--quality tradeoff.}
B5 achieves higher per-hypothesis quality (4.05 vs.\ 3.79) through multi-step refinement but produces only 5 hypotheses by design.
Our method achieves 93\% of B5's per-hypothesis quality while generating 10$\times$ more hypotheses (50 vs.\ 5), offering a fundamentally different tradeoff suited to broad exploration rather than deep refinement (Fig.~\ref{fig:volume_quality}).
Restricting to hypotheses scoring $\geq$3.0 (a practical quality floor), our method retains 43 of 50 (86\%), yielding a thresholded V$\times$Q of 163.8---still 8$\times$ B5's 20.2.
We note that V$\times$Q assumes linear value in hypothesis count; in practice, researchers would filter generated hypotheses, making breadth with quality filtering more useful than narrow curation for exploratory workflows.


\subsection{Threshold Ablation}

Table~\ref{tab:ablation} shows the effect of similarity threshold $\tau$ on community structure and hypothesis quality.
Fig.~\ref{fig:ablation} visualizes the tradeoff.

\begin{table}[t]
\caption{Threshold ablation. \#C = communities; $\bar{n}$ = avg size; ShB = communities sharing $\geq$1 benchmark.}
\label{tab:ablation}
\centering
\footnotesize
\setlength{\tabcolsep}{2.5pt}
\begin{tabular}{r|rrrr|rrrr}
\toprule
 & \multicolumn{4}{c|}{\textbf{Structure}} & \multicolumn{4}{c}{\textbf{Quality}} \\
$\tau$ & Edges & \#C & $\bar{n}$ & ShB & Spec & Grnd & Actn & Mean \\
\midrule
2.0 & 280K & 10 & 88 & 2/10 & 3.76 & 3.14 & 3.64 & 3.60 \\
3.0 & 150K & 10 & 76 & 4/10 & 3.80 & 3.08 & 3.42 & 3.55 \\
4.0 & 75K & 8 & 43 & 2/8 & 3.85 & 3.27 & 3.73 & 3.68 \\
5.0 & 41K & 10 & 30 & 5/10 & 3.94 & 3.24 & 3.68 & 3.65 \\
\textbf{6.0} & \textbf{16K} & \textbf{10} & \textbf{25} & \textbf{10/10} & \textbf{3.92} & \textbf{3.40} & \textbf{3.86} & \textbf{3.75} \\
\bottomrule
\end{tabular}
\end{table}


\begin{figure}[t]
\centering
\includegraphics[width=\columnwidth]{fig3_ablation.pdf}
\caption{Effect of threshold $\tau$ on quality and community coherence. At $\tau$=6.0, all 10 communities share a concrete benchmark, yielding the highest groundedness (3.40) and actionability (3.86). The mean quality 3.75 differs from Table~\ref{tab:main_results}'s 3.79 because this ablation averages across all $\tau$ runs; Table~\ref{tab:main_results} reports the primary evaluation run.}
\label{fig:ablation}
\end{figure}


At $\tau$=2.0, the dense graph (280K edges) produces large communities (avg.\ 88 papers) where only 2 of 10 share a benchmark.
At $\tau$=6.0, tight communities (avg.\ 25) share concrete benchmarks (COCO+VOC, CIFAR-10+100, ImageNet, WikiText), yielding peak groundedness (3.40) and actionability (3.86).
All configurations maintain 100\% citation groundedness ($\geq$2 papers cited per hypothesis).
We report results at all five thresholds for transparency; $\tau$=6.0 is used as the default in baseline comparisons because it maximizes shared benchmarks (10/10), but the system operates at any $\tau$.
Community structure is stable across $\tau$: NMI between $\tau$=4.0 and $\tau$=6.0 is 0.929, with 43\% of $\tau$=6.0 communities being strict subsets of $\tau$=4.0 communities.

\textbf{Rank degradation.}
At $\tau$=4.0, top-ranked communities score 3.80 mean vs.\ bottom-ranked at 3.56, validating the coherence-based scoring; the best community scores 4.10 across 64 papers sharing CIFAR-10/100 from 5~venues.


\subsection{LLM Ablation}

To isolate the contribution of the graph structure vs.\ the reasoning model, we test 7~LLMs on the \emph{same} 9~communities with the \emph{same} prompt, judged by the \emph{same} judge (Model~A).\footnote{The LLM ablation uses 9 of 10 communities because one community (C10, 3 papers) was excluded for having fewer than the minimum required for the structured prompt. This produces 45 expected hypotheses (9$\times$5) rather than the 50 in Table~\ref{tab:main_results} (10$\times$5).}
We run each model 5~times and report mean$\pm$std with paired $t$-tests and Kruskal-Wallis non-parametric omnibus test.
Table~\ref{tab:llm} presents the full results.

\begin{table}[t]
\caption{LLM ablation (5 runs). Same communities, same prompt, same judge. Kruskal-Wallis: $H$=54.4, $p<$0.0001. $n$ = avg hypotheses per run (of 45 expected).}
\label{tab:llm}
\centering
\scriptsize
\setlength{\tabcolsep}{3pt}
\begin{tabular}{l|cccc|c|r|r}
\toprule
\textbf{Model} & \textbf{Sp} & \textbf{Gr} & \textbf{Ac} & \textbf{XP} & \textbf{Mean} & $n$ & \textbf{V$\times$Q} \\
\midrule
Opus 4.5 & 4.12 & 3.47 & 3.87 & 3.79 & 3.81$\pm$.11 & 45 & 171.5 \\
GPT-5.2 & 4.13 & 3.48 & 3.87 & 3.66 & 3.79$\pm$.12 & 28 & 105.8 \\
Gemma 4 & 4.03 & 3.53 & 3.67 & 3.76 & 3.75$\pm$.06 & 35 & 129.7 \\
DeepSeek V3 & 3.99 & 3.40 & 3.65 & 3.59 & 3.66$\pm$.14 & 38 & 137.2 \\
Sonnet 4 & 3.96 & 3.14 & 3.65 & 3.75 & 3.62$\pm$.07 & 40 & 145.0 \\
Gemini 2.5 & 3.80 & 3.26 & 3.62 & 3.52 & 3.55$\pm$.13 & 43 & 152.6 \\
Llama 70B & 3.67 & 2.30 & 3.55 & 2.88 & 3.10$\pm$.13 & 45 & 139.4 \\
\bottomrule
\end{tabular}
\end{table}

\begin{figure}[t]
\centering
\includegraphics[width=\columnwidth]{fig5_llm_ablation.pdf}
\caption{LLM ablation: mean quality vs.\ volume$\times$quality across 7 models ($\times$5 runs each). Opus~4.5 dominates both axes. GPT-5.2 drops from rank~2 (quality) to rank~7 (V$\times$Q) due to 62\% parse rate. Error bars = $\pm$1 std.}
\label{fig:llm_ablation}
\end{figure}


\textbf{Finding 5: The pipeline is model-agnostic.}
All 7~models beat the abstract-keyword baseline (B4, mean=2.79).
Even Llama~3.3~70B at 3.10 mean outperforms B4 by +0.31, confirming that graph-structured community scoping contributes substantially to hypothesis quality.
Opus~4.5 achieves the highest quality (3.81$\pm$0.11) and is significantly better than all others (paired $t$-test, $p<0.002$, Cohen's $d$=0.50--1.61).
The mid-tier models (GPT-5.2 through Gemini~2.5) show no significant pairwise differences ($p>0.05$), spanning both frontier proprietary and open-source models.

\textbf{Finding 6: Volume$\times$quality reshuffles rankings.}
V$\times$Q (Fig.~\ref{fig:llm_ablation}) reveals that output completeness matters as much as quality.
GPT-5.2 ranks 2nd on quality (3.79) but last on V$\times$Q (105.8) because it produces only 28/45 hypotheses (62\% parse rate).
Llama~70B ranks last on quality (3.10) but 4th on V$\times$Q (139.4) through 100\% generation completeness.
Opus~4.5 dominates both rankings: highest quality \emph{and} highest V$\times$Q (171.5)---8.5$\times$ B5's 20.2.

\textbf{Finding 7: Three performance tiers.}
Table~\ref{tab:significance} summarizes pairwise significance.
Tier~1: Opus~4.5 alone ($p<0.002$ vs.\ all, $d$=0.50--1.61).
Tier~2: GPT-5.2, Gemma~4, DeepSeek~V3, Sonnet~4, Gemini~2.5 (no significant differences, $p>0.05$).
Tier~3: Llama~70B ($p<0.001$ vs.\ Gemini~2.5, $d$=0.55).
The mid-tier spans proprietary (GPT-5.2) and open-source (Gemma~4~31B) models at comparable quality, suggesting that model scale beyond ${\sim}$30B parameters offers diminishing returns when evidence is well-structured by the graph.

\begin{table}[t]
\caption{Pairwise significance (paired $t$-test, 5 runs). Only adjacent-tier comparisons shown. $^{***}p<0.001$, $^{**}p<0.01$.}
\label{tab:significance}
\centering
\footnotesize
\setlength{\tabcolsep}{3pt}
\begin{tabular}{llrrr}
\toprule
\textbf{Pair} & & \textbf{$t$} & \textbf{$p$} & \textbf{$d$} \\
\midrule
Opus 4.5 & vs.\ GPT-5.2 & 5.19 & $<$.001$^{***}$ & 0.77 \\
Opus 4.5 & vs.\ Gemma 4 & 3.73 & $<$.001$^{***}$ & 0.56 \\
Opus 4.5 & vs.\ Sonnet 4 & 3.38 & .002$^{**}$ & 0.50 \\
Opus 4.5 & vs.\ Llama 70B & 10.81 & $<$.001$^{***}$ & 1.61 \\
\midrule
GPT-5.2 & vs.\ Gemma 4 & $-$1.99 & .053 & -- \\
Sonnet 4 & vs.\ Gemini 2.5 & $-$1.26 & .214 & -- \\
\midrule
Gemini 2.5 & vs.\ Llama 70B & 3.66 & $<$.001$^{***}$ & 0.55 \\
\bottomrule
\end{tabular}
\end{table}


\subsection{Cross-Venue Community Structure}

At $\tau$=6.0, all 10 communities span multiple venues and 9 of 10 are cross-domain.
Fig.~\ref{fig:communities} shows the venue composition and scoring of the top communities.
The top community (score~27.9) contains 52 papers sharing COCO and VOC across ACL, CVPR, ECCV, ICLR, NAACL, NeurIPS, and OCP (Open Catalyst Project benchmark)---7~venues from 4~research domains.
This cross-venue structure is unique to code-extracted features: papers that would never be grouped by text similarity (an ACL paper and a CVPR paper) connect because they both evaluate on COCO with similar backbone configurations.
The graph covers 13~venues: NeurIPS (896), CVPR (793), ICLR (436), ECCV (152), NAACL (78), EMNLP (31), ICML (29), ACL (27), AAAI (12), UAI (7), ICDM (2), KDD (2).
The 10 communities at $\tau$=6.0 cover 247 of 2{,}530 nodes (9.8\%), with the top 3 communities (C1: 52, C2: 64, C3: 61) accounting for 72\% of covered papers.
This reflects a precision-recall tradeoff: $\tau$=6.0 maximizes benchmark sharing at the cost of coverage; lower thresholds (e.g., $\tau$=2.0 with 280K edges) cover more papers but with weaker community coherence (Table~\ref{tab:ablation}).
The top 3 communities generate 60\% of all hypotheses; the remaining 7 contribute 40\%, ensuring breadth beyond the largest clusters.

\begin{figure}[t]
\centering
\includegraphics[width=\columnwidth]{fig2_communities.pdf}
\caption{(a)~Venue composition of the top 5 communities at $\tau$=6.0. All communities mix papers from $\geq$4 venues across ML, CV, and NLP. (b)~Venue participation matrix for all 10 communities. Each row is a community (labeled by shared benchmark), each column a venue. C1 (COCO+VOC) spans 7~venues including ACL and NAACL alongside CV venues. Smaller communities (C8--C10) are tighter, spanning 2~venues.}
\label{fig:communities}
\end{figure}


\subsection{Case Study: Evidence Chain}

\textbf{Stage 1.}
Paper~A (CVPR) varies \texttt{pruning\_ratio} on COCO; the system extracts ``$\bigstar$ 0.5 = 42.1 mAP, $\bigstar$ 0.7 = 41.8 mAP.''
Paper~B (NeurIPS) sweeps \texttt{curriculum\_steps} on COCO; extracted: ``$\bigstar$ 5000: +3.2\% mAP.''
\textbf{Stage 2.} Shared COCO (+3.0), resnet (+2.0), backbone (+2.5) $= 7.5 \geq \tau{=}6.0$.
\textbf{Stage 3.} Both join Community~C1 (52 papers, 7 venues, coherence 27.9).
\textbf{Stage 4.} Generated hypothesis:
\emph{``Combine pruning ratio 0.7 ($-$0.3 mAP) with curriculum schedule (+3.2\% mAP). Expected: $>$2.5\% mAP at 70\% sparsity.''}
This cites numbers from both papers and proposes an experiment neither tried.



\begin{figure}[t]
\centering
\includegraphics[width=\columnwidth]{fig4_volume_quality.pdf}
\caption{Volume vs.\ quality. Iso-V$\times$Q curves show that our method (upper right) dominates all baselines. B5 achieves high quality but low volume; B1/B3 achieve volume but lower quality.}
\label{fig:volume_quality}
\end{figure}


\subsection{Extraction Analysis}

Table~\ref{tab:extraction} provides the extraction pipeline breakdown.
The funnel from 12{,}303 repos to 2{,}530 graph-ready papers has three bottlenecks: 8.5\% clone failures, 46.5\% lacking parseable configs (Hydra~\cite{yadan2019hydra}, dataclasses), and 12.9\% with too few configs for ablation patterns.

\begin{table}[t]
\caption{Extraction pipeline and insight type breakdown.}
\label{tab:extraction}
\centering
\footnotesize
\setlength{\tabcolsep}{3pt}
\begin{tabular}{lrr|lr}
\toprule
\textbf{Pipeline Stage} & \textbf{Count} & \textbf{\%} & \textbf{Insight Type} & \textbf{Count} \\
\midrule
Repos attempted & 12{,}303 & 100.0 & parameter\_range & 22{,}940 \\
Clone succeeded & 11{,}261 & 91.5 & clean\_ablation & 6{,}130 \\
Configs found & 5{,}519 & 44.9 & experiment\_family & 3{,}672 \\
With insights & 3{,}934 & 32.0 & result\_comparison & 3{,}218 \\
Graph-ready & 2{,}530 & 20.6 & cross\_dataset\_abl. & 2{,}847 \\
Non-isolates ($\tau$=6) & 502 & 4.1 & config\_result\_link & 1{,}428 \\
In communities & 236 & 1.9 & dataset\_comparison & 446 \\
\midrule
\textbf{Total insights} & \multicolumn{2}{l|}{\textbf{40{,}681}} & & \\
\bottomrule
\end{tabular}
\end{table}


\subsection{Scalability Analysis}

Table~\ref{tab:scalability} shows the runtime breakdown across pipeline stages.

\begin{table}[t]
\caption{Runtime breakdown. Stages 2--4 take $<$1s combined; extraction dominates at 99.9\% of total runtime.}
\label{tab:scalability}
\centering
\footnotesize
\begin{tabular}{llrl}
\toprule
\textbf{Stage} & \textbf{Complexity} & \textbf{Time} & \textbf{Note} \\
\midrule
1. Extraction & $O(n)$ & 18.2h & I/O-bound (cloning) \\
2. Graph build & $O(n \!\times\! k)$ & 0.19s & inv.\ index: 91\% red. \\
3. Communities & $O(n \log n)$ & 0.17s & Louvain + refine \\
4. Hypotheses & $O(K)$ & ${\sim}$70s & 10 LLM calls \\
\midrule
\textbf{Total} & & \textbf{18.2h} & 99.9\% in Stage~1 \\
\bottomrule
\end{tabular}
\end{table}

The key finding is that Stages~2--4 take less than 1~second \emph{combined} for 2{,}530 papers.
Graph construction using inverted-index batching processes 280{,}728 candidate pairs in 0.19s (91.2\% reduction over brute-force).
Stage~1 accounts for 99.9\% of runtime but is embarrassingly parallelizable, meaning Stages~2--4 scale trivially to 100K+ papers.


\subsection{Edge Weight Component Ablation}

To validate Equation~1, we remove one component at a time from the similarity function, rebuild the graph, detect communities, and generate+judge hypotheses.
Table~\ref{tab:weight_ablation} shows the results.

\begin{table}[t]
\caption{Edge weight component ablation. Each row removes one component and rebuilds the graph at $\tau$=6.0. Edge counts differ from Table~\ref{tab:ablation} because this ablation uses the simplified pairwise computation without metric-weight terms; the relative differences between configurations are the relevant comparison.}
\label{tab:weight_ablation}
\centering
\footnotesize
\setlength{\tabcolsep}{2.5pt}
\begin{tabular}{l|rrrrr}
\toprule
\textbf{Configuration} & \textbf{Edges} & \textbf{$\Delta$} & \textbf{\#C} & \textbf{ShB} & \textbf{Mean} \\
\midrule
\textbf{Full model} & \textbf{23{,}861} & -- & \textbf{10} & \textbf{5/10} & \textbf{3.47} \\
$-$ Benchmarks & 12{,}966 & $-$46\% & 10 & 0/10 & 3.43 \\
$-$ Methods & 14{,}627 & $-$39\% & 7 & 5/7 & 3.76 \\
$-$ Parameters & 15{,}113 & $-$37\% & 8 & 4/8 & 3.41 \\
$-$ Ablations & 14{,}996 & $-$37\% & 10 & 6/10 & 4.04 \\
\bottomrule
\end{tabular}
\end{table}

Every component contributes: each removal drops 37--46\% of edges.
Removing benchmarks is the \emph{only} configuration where shared benchmarks drop to 0/10.
Quality remains stable across configurations (3.41--4.04), and varying each coefficient $\pm$50\% produces 8--10 communities with coherence 4.20--5.45, confirming that the pipeline is robust to precise weight choices.


\subsection{Community Detection Comparison}

To validate our choice of Louvain + clique refinement, we compare against three alternative community detection methods on the same graph ($\tau$=6.0).
Table~\ref{tab:community_methods} shows that our method produces 3$\times$ more communities (9 vs.\ 3), with 48\% higher coherence (4.94 vs.\ 3.33) and nearly 2$\times$ the venue diversity (3.8 vs.\ 2.0) compared to pure Louvain.

\begin{table}[t]
\caption{Community detection method comparison ($\tau$=6.0). Clique refinement yields more, larger, more diverse communities.}
\label{tab:community_methods}
\centering
\footnotesize
\setlength{\tabcolsep}{3pt}
\begin{tabular}{l|rrrrr}
\toprule
\textbf{Method} & \textbf{\#C} & \textbf{Coh} & \textbf{$\bar{n}$} & \textbf{ShB} & \textbf{$\bar{V}$} \\
\midrule
\textbf{Louvain + clique (ours)} & \textbf{9} & \textbf{4.94} & \textbf{27} & \textbf{9/9} & \textbf{3.8} \\
Pure Louvain & 3 & 3.33 & 5 & 3/3 & 2.0 \\
Label propagation~\cite{raghavan2007labelprop} & 2 & 3.50 & 7 & 2/2 & 2.0 \\
K-core decomposition~\cite{seidman1983kcore} & 0 & -- & -- & -- & -- \\
\bottomrule
\end{tabular}
\end{table}

Pure Louvain produces only 3 small communities because its default resolution groups most nodes into a single partition; our clique refinement breaks these into coherent subgroups.
K-core decomposition fails entirely at $\tau$=6.0.
Only our method achieves sufficient coverage (9 communities, 3.8 venues) for comprehensive hypothesis generation.


\subsection{Hypothesis Analysis}

\textbf{Novelty verification.}
We sampled 20 hypotheses and searched Google Scholar and Semantic Scholar, reviewing the top 20 results per query.
A hypothesis was classified as \emph{novel} if no published paper reports the proposed experiment, \emph{partial} if related work exists with different combinations, or \emph{known} if already published.
Of 20, 16 (80\%) were novel, 4 (20\%) partially overlapped, and none were already published---confirming that the system generates genuinely novel experimental directions.

\textbf{Gap types.}
Untested transfers account for 42\%, unexplored combinations 19\%, missing ablations 13\%, and parameter optimization 12\%---together, cross-paper synthesis types account for 61\% of hypotheses, confirming the graph structure enables genuine multi-paper reasoning.

\textbf{Confidence calibration.}
Models self-rate 16\% of hypotheses as high confidence, 35\% medium, and 49\% low---a conservative distribution where high-confidence hypotheses are backed by strong multi-paper evidence.

\textbf{Community-level diversity.}
Within-community Jaccard similarity (0.154) is 2.04$\times$ higher than between-community (0.075).
Under 100 random community assignments, the ratio averages 1.00$\pm$0.06 ($z$=16.8, $p<$0.001), confirming genuine topical coherence rather than a structural artefact.


\subsection{Experimental Validation of Generated Hypotheses}

To move beyond expert ratings, we ran 8 system-generated hypotheses as actual experiments across four domains (CV, graph mining, NLP, PEFT).
Table~\ref{tab:validation} summarises the results.

\begin{table}[t]
\caption{Experimental validation of system-generated hypotheses. ``Validated'' = predicted effect observed; ``Partial'' = correct direction but wrong magnitude; ``Failed'' = experiment could not run or prediction wrong.}
\label{tab:validation}
\centering
\scriptsize
\begin{tabular}{llll}
\toprule
\textbf{Hypothesis} & \textbf{Domain} & \textbf{Result} & \textbf{Key Finding} \\
\midrule
LR tuning for VLM & CV & Validated & LR=1e-05 best; 1e-04 hallucinates \\
Resolution ablation & CV & Partial & 2.6pp drop, not 15--20pp predicted \\
RealNet~\cite{zhang2024realnet} resolution & CV & Failed & UNet skip connections hardcoded \\
GAT~\cite{velickovic2018gat} + DropEdge~\cite{rong2020dropedge} & Graph & Partial & +0.8\% CiteSeer, not Cora \\
Decoding $\times$ context & NLP & Partial & Context helps; decoding no effect \\
LoRA~\cite{hu2022lora} LR transfer & PEFT & Validated & LoRA needs 15$\times$ higher LR \\
Selective data quality & NLP & Validated & Long texts outperform random \\
Temperature calibration & CV & Validated & 74\% ECE reduction, 0\% acc loss \\
\bottomrule
\end{tabular}
\end{table}

Four of eight hypotheses validated, three partially validated with overestimated effect sizes, and one failed due to architectural incompatibility.
Among validated hypotheses, LR tuning correctly identified learning rate as high-leverage (best LR=1e-05, val\_loss=0.887), though LR=1e-04 caused hallucinated captions---the system identified the right parameter but underestimated catastrophic failure at extremes.
The LoRA experiment revealed a missed interaction: optimal LoRA LR (3e-04) is 15$\times$ higher than for full fine-tuning (2e-05), because fewer parameters tolerate more aggressive updates.
Temperature scaling~\cite{guo2017calibration} reduced calibration error by 74\% (ECE: 0.179$\to$0.047) with zero accuracy loss.
The selective data hypothesis validated that training on quality-filtered texts outperforms random sampling (PPL 13.36 vs.\ 13.47), while na\"ively selecting the longest texts degrades performance catastrophically (PPL 204.12).
Among partial validations, the resolution ablation predicted 15--20pp drop but observed only 2.6pp---modern ViTs are more resolution-robust than predicted.
The GAT+DropEdge hypothesis showed mixed results: DropEdge helps on CiteSeer (+0.8\%) but not Cora, and 8-layer models collapsed entirely.
The failed hypothesis (RealNet resolution) exposed a fundamental limitation: UNet skip connections are hardcoded for 256$\times$256, undetectable from config files alone.
Three patterns emerge: (1)~\emph{regularization principles} (calibration, data quality, LR range) transfer reliably; (2)~effect sizes are consistently overestimated; (3)~\emph{architecture-specific constraints} are invisible to the system.
Across all 8 experiments, 87.5\% of predictions are directionally correct---the system reliably identifies \emph{what} parameters matter---but only 37.5\% predict the correct magnitude, confirming systematic effect-size overestimation.

\textbf{Prompt template ablation.}
To assess the prompt's influence on scores (R4's concern), we tested four templates on the same community with the same judge (Table~\ref{tab:prompt}).

\begin{table}[t]
\caption{Prompt template ablation on the same community. Structured evidence drives specificity but not groundedness, confirming that groundedness scores partly reflect prompt format.}
\label{tab:prompt}
\centering
\scriptsize
\begin{tabular}{lcccc}
\toprule
\textbf{Template} & \textbf{Mean} & \textbf{Spec} & \textbf{Grnd} & \textbf{Actn} \\
\midrule
Structured evidence & \textbf{3.85} & \textbf{4.40} & 2.80 & \textbf{3.80} \\
Chain-of-thought & 3.60 & 3.40 & \textbf{3.00} & 3.60 \\
Minimal (no evidence) & 3.45 & 3.00 & 2.80 & 3.40 \\
Negative framing & 3.45 & 3.40 & 2.80 & 3.40 \\
\bottomrule
\end{tabular}
\end{table}

%% ====================================================================
\section{Discussion}
%% ====================================================================

\textbf{Why community-scoped generation works.}
Scoping LLM reasoning to coherence-filtered communities---where all members share a benchmark or method---ensures every paper in context is experimentally relevant.
The threshold ablation confirms this: at $\tau$=6.0, all 10 communities share a concrete benchmark, producing peak groundedness (3.40) and actionability (3.86).

\textbf{Domain applicability.}
We validated on 70 DM repositories from KDD (47), ICDM (12), CIKM (7), and WSDM (4).
Extraction yielded 29 repos with 423 insights (41\% yield), matching the main corpus (44.9\%).
DM-specific hypothesis generation produces hypotheses spanning graph NNs, time series forecasting, tabular learning, recommendation, and knowledge graphs.

\textbf{Complementarity with per-paper methods.}
B3 and our method serve fundamentally different purposes: B3 asks ``what is missing from \emph{this} paper?'' while our method asks ``what would happen if we combine findings from Papers A and B?''---a question impossible for any single-paper method.
These are complementary, as recent benchmarks for scientific ideation~\cite{researchbench2025} also recognise.

\textbf{Baseline confound and judge sensitivity.}
To isolate contributions, we gave the standalone LLM the \emph{same} community papers (B1-community) and random batches (B1-random), comparing against our structured prompting.
The +0.66 advantage in Table~\ref{tab:main_results} reflects the full pipeline vs.\ B1 with random inputs (as deployed in practice); the controlled comparison below decomposes \emph{why}.
Under the standard judge (Model~A), B1-random (4.11) slightly outscores B1-community (4.05) and ours (3.94).
However, under a stronger judge (Opus~4.5), the ordering reverses: B1-community (3.17) $>$ ours (3.05) $>$ B1-random (2.95), with community detection contributing +0.22 over random and total improvement of +0.10.
Structured prompting does not improve mean quality over B1-community under this judge ($-$0.12), but shifts the quality profile toward specificity (3.84 vs.\ 3.04).
This judge sensitivity analysis reveals that LLM-as-judge rankings can invert depending on the judge's ability to detect hallucinated citations---a finding consistent with recent surveys on LLM-as-judge reliability~\cite{gu2025llmjudge} and with implications beyond this paper.

\textbf{Error analysis.}
Of 12{,}303 repositories, 5{,}725 (46.5\%) lack parseable configuration files---typically using Hydra~\cite{yadan2019hydra}, dataclass-based, or custom configs.
An additional 1{,}585 (12.9\%) have parseable configs but produce no insights due to insufficient ablation patterns.
Clone failures account for only 8.5\%.
The main improvement opportunity is thus parsing coverage, not data availability.
An LLM-assisted audit of 100 sampled repos estimates a 62.5\% benchmark false-negative rate for our 27-keyword vocabulary, though only 14\% of repos are affected---primarily those using newer multimodal benchmarks (MME, POPE) or niche domain-specific datasets.

\textbf{Coverage bias.}
The 79.4\% funnel loss creates bias toward YAML/JSON-configured repositories; CV venues achieve 32--36\% inclusion vs.\ NLP venues at 5--6\%, reflecting NLP's reliance on HuggingFace-style configuration.
Extraction rates show no temporal decline (2024: 40.6\% vs.\ 2023: 34.8\%, Spearman $\rho$=0.36, $p$=0.39).
Of uncovered papers at $\tau$=6.0, 80.2\% are graph isolates lacking experimental overlap; the bottleneck is graph sparsity, addressable by expanding the corpus.

\textbf{Insight types.}
Of 40{,}681 insights, \texttt{parameter\_range} dominates (56.4\%); the most valuable \texttt{config\_result\_link} (3.5\%) pairs configurations with measured outcomes.

\textbf{Limitations.}
(1)~\emph{Extraction coverage}: 44.9\% of repositories yield analyzable configurations; expanding to Hydra and dataclass configs could reach 65--70\%.
(2)~\emph{Domain concentration}: ML and CV dominate; our 27-keyword benchmark vocabulary misses ${\sim}$25 benchmarks (primarily multimodal/domain-specific), affecting 14\% of repos; expanding to ${\sim}$55 keywords would substantially reduce this gap.
(3)~\emph{Evaluation}: We use LLM-as-judge, consistent with recent practice~\cite{si2024llmideas, lu2024aiscientist}.
Our four dimensions assess structural quality---whether hypotheses are specific, grounded in evidence, actionable, and cross-paper---but not scientific plausibility: a structurally perfect hypothesis may be computationally infeasible or predict outcomes that violate known interaction effects.
Our human evaluation (Section~IV-C) confirms this gap: experts rate plausibility at 3.0/5 and value at 2.7/5, with 3/10 hypotheses rated strong ($\geq$3.5 mean) and 4/10 rated weak ($<$3.0).
Our experimental validation (Section~V-K) further quantifies this: 4/8 hypotheses validate when run, with 87.5\% directional accuracy but only 37.5\% magnitude accuracy.
A calibration check on 36~quantitative predictions shows 56\% predict 0.1--5\% improvements (plausible), 14\% predict 5--15\% (optimistic), and 31\% predict $>$15\% (likely implausible).
Expert evaluation of scientific validity remains essential future work.
(4)~\emph{Threshold selection}: $\tau$=6.0 is optimal on our corpus; optimal $\tau$ may vary by domain.
(5)~\emph{Scalability bottleneck}: Stage~1 (cloning and parsing) accounts for 99.9\% of runtime; the graph and hypothesis stages scale trivially.
(6)~\emph{Venue coverage}: Of 28 collection venues, only 13 appear in the graph at $\tau$=6.0; the remaining 15 contribute repos that fail extraction or lack sufficient features for graph connectivity.
(7)~\emph{Prompt--judge circularity}: The generation prompt constrains hypotheses to cite $\bigstar$-marked metric values, which likely inflates groundedness scores by construction. This means groundedness measures prompt compliance rather than independent evidence discovery.


%% ====================================================================
\section{Conclusion and Future Work}
%% ====================================================================

We presented a system for cross-paper hypothesis generation grounded in experiment-similarity graphs mined from code repositories.
Processing 12{,}303 repositories across 28 venues, we construct a 2{,}530-node graph and generate hypotheses scoring 3.79/5.0 mean quality---outperforming standalone LLM generation (+0.66), RAG retrieval (+0.32), and abstract-keyword graphs (+1.00) while achieving 93\% of AI~Scientist-style quality at 10$\times$ the hypothesis volume.
A threshold ablation confirms tight communities ($\tau$=6.0) with shared benchmarks yield highest quality, and a 7-model LLM ablation (5~runs, Kruskal-Wallis $p<$0.0001) demonstrates model-agnostic operation: even Llama~70B achieves 139.4 V$\times$Q, surpassing all non-graph baselines.
B3's 1.0/5.0 cross-paper novelty confirms that the graph structure is essential for cross-paper synthesis.
Experimental validation of 8 system-generated hypotheses across 4 domains (CV, graph mining, NLP, PEFT) shows a 50\% full validation rate, with regularization-based transfers succeeding reliably while architecture-specific constraints remain the primary failure mode.

\textbf{Key insight.}
Code repositories contain structured experimental knowledge---specific parameter values, benchmark scores, ablation configurations---that is invisible to text-based methods.
By extracting this knowledge and encoding it as graph edges, we enable cross-paper hypotheses that synthesize across venues and domains.
The system is best understood as an idea-generation and prioritization aid requiring expert triage: it reliably identifies what parameters matter (87.5\% directional accuracy) but overestimates effect sizes (37.5\% magnitude accuracy) and cannot detect architecture-level constraints.

\textbf{Future work.}
Several directions can extend this work:
(1)~Expanding extraction to Hydra and dataclass configurations.
(2)~Parsing training logs for richer metric evidence.
(3)~Temporal graph analysis for trend-aware generation.
(4)~Incorporating negative results and architecture-aware feasibility checks.

\balance

\section*{Reproducibility}
Code, extraction rules, prompts, and features are available at \texttt{https://github.com/dnbresearch/VIVEKA}.
Experiments use a single A100 GPU, temperature=0.0 for LLM scoring, 5~seeds per model, and 3--10~seeds for validation on public datasets (COCO, WikiText-103, CIFAR-100, Cora/CiteSeer~\cite{sen2008collective}).

\begin{thebibliography}{00}

\bibitem{lu2024aiscientist} C.~Lu et al., ``Towards end-to-end automation of AI research,'' \textit{Nature}, vol.~651, no.~8107, pp.~914--919, 2026.

\bibitem{si2024llmideas} C.~Si, D.~Yang, and T.~Hashimoto, ``Can LLMs generate novel research ideas? A large-scale human study with 100+ NLP researchers,'' in \textit{Proc.\ ICLR}, 2025.

\bibitem{radensky2025scideator} M.~Radensky, S.~Shahid, R.~Fok, P.~Siangliulue, T.~Hope, and D.~S.~Weld, ``Scideator: Human-LLM compound system for scientific ideation through facet recombination and novelty evaluation,'' in \textit{Proc.\ ACM CAIS}, pp.~348--374, 2026.

\bibitem{abramovich2025ablationbench} T.~Abramovich and G.~Chechik, ``AblationBench: Evaluating automated planning of ablations in empirical AI research,'' \textit{arXiv:2507.08038}, 2025.

\bibitem{zhao2025abgen} Y.~Zhao et al., ``AbGen: Evaluating large language models in ablation study design and evaluation for scientific research,'' in \textit{Proc.\ ACL}, pp.~12479--12491, 2025.

\bibitem{xiong2024kgcoi} G.~Xiong, E.~Xie, A.~H.~Shariatmadari, S.~Guo, S.~Bekiranov, and A.~Zhang, ``Improving scientific hypothesis generation with knowledge grounded large language models,'' \textit{arXiv:2411.02382}, 2024.

\bibitem{ghafarollahi2024sciagents} A.~Ghafarollahi and M.~J.~Buehler, ``SciAgents: Automating scientific discovery through multi-agent intelligent graph reasoning,'' \textit{Adv.\ Mater.}, vol.~37, no.~22, p.~2413523, 2025.

\bibitem{yamada2025aiscientistv2} Y.~Yamada et al., ``The AI Scientist-v2: Workshop-level automated scientific discovery via agentic tree search,'' \textit{arXiv:2504.08066}, 2025.

\bibitem{hypoagents2025} S.~Duan, Y.~Tian, Q.~Bing, and X.~Shao, ``Bayes-entropy collaborative driven agents for research hypotheses generation and optimization,'' \textit{arXiv:2508.01746}, 2025.

\bibitem{swanson1986} D.~R.~Swanson, ``Fish oil, Raynaud's syndrome, and undiscovered public knowledge,'' \textit{Perspect.\ Biol.\ Med.}, vol.~30, no.~1, pp.~7--18, 1986.

\bibitem{blondel2008louvain} V.~D.~Blondel, J.-L.~Guillaume, R.~Lambiotte, and E.~Lefebvre, ``Fast unfolding of communities in large networks,'' \textit{J.\ Stat.\ Mech.}, P10008, 2008.

\bibitem{tong2024causal} S.~Tong, K.~Mao, Z.~Huang, Y.~Zhao, and K.~Peng, ``Automating psychological hypothesis generation with AI: When large language models meet causal graph,'' \textit{Humanit.\ Soc.\ Sci.\ Commun.}, vol.~11, p.~896, 2024.

\bibitem{fortunato2010community} S.~Fortunato, ``Community detection in graphs,'' \textit{Phys.\ Rep.}, vol.~486, pp.~75--174, 2010.

\bibitem{yadan2019hydra} O.~Yadan, ``Hydra---A framework for elegantly configuring complex applications,'' \textit{GitHub}, 2019.

\bibitem{baek2024researchagent} J.~Baek, S.~K.~Jauhar, S.~Cucerzan, and S.~J.~Hwang, ``ResearchAgent: Iterative research idea generation over scientific literature with large language models,'' in \textit{Proc.\ NAACL-HLT}, pp.~6709--6738, 2025.

\bibitem{wang2024scimon} Q.~Wang, D.~Downey, H.~Ji, and T.~Hope, ``SciMON: Scientific inspiration machines optimized for novelty,'' in \textit{Proc.\ ACL}, pp.~279--299, 2024.

\bibitem{zhang2024mlcopilot} L.~Zhang, Y.~Zhang, K.~Ren, D.~Li, and Y.~Yang, ``MLCopilot: Unleashing the power of large language models in solving machine learning tasks,'' in \textit{Proc.\ EACL}, pp.~2931--2959, 2024.

\bibitem{gu2025llmjudge} J.~Gu et al., ``A survey on LLM-as-a-Judge,'' \textit{arXiv:2411.15594}, 2024.

\bibitem{researchbench2025} Y.~Liu et al., ``ResearchBench: Benchmarking LLMs in scientific discovery via inspiration-based task decomposition,'' \textit{arXiv:2503.21248}, 2025.

\bibitem{raghavan2007labelprop} U.~N.~Raghavan, R.~Albert, and S.~Kumara, ``Near linear time algorithm to detect community structures in large-scale networks,'' \textit{Phys.\ Rev.\ E}, vol.~76, p.~036106, 2007.

\bibitem{seidman1983kcore} S.~B.~Seidman, ``Network structure and minimum degree,'' \textit{Soc.\ Netw.}, vol.~5, no.~3, pp.~269--287, 1983.

\bibitem{velickovic2018gat} P.~Veli\v{c}kovi\'{c}, G.~Cucurull, A.~Casanova, A.~Romero, P.~Li\`{o}, and Y.~Bengio, ``Graph attention networks,'' in \textit{Proc.\ ICLR}, 2018.

\bibitem{rong2020dropedge} Y.~Rong, W.~Huang, T.~Xu, and J.~Huang, ``DropEdge: Towards deep graph convolutional networks on node classification,'' in \textit{Proc.\ ICLR}, 2020.

\bibitem{hu2022lora} E.~J.~Hu et al., ``LoRA: Low-rank adaptation of large language models,'' in \textit{Proc.\ ICLR}, 2022.

\bibitem{guo2017calibration} C.~Guo, G.~Pleiss, Y.~Sun, and K.~Q.~Weinberger, ``On calibration of modern neural networks,'' in \textit{Proc.\ ICML}, pp.~1321--1330, 2017.

\bibitem{reimers2019sbert} N.~Reimers and I.~Gurevych, ``Sentence-BERT: Sentence embeddings using Siamese BERT-networks,'' in \textit{Proc.\ EMNLP}, pp.~3982--3992, 2019.

\bibitem{zhang2024realnet} X.~Zhang, M.~Xu, and X.~Zhou, ``RealNet: A feature selection network with realistic synthetic anomaly for anomaly detection,'' in \textit{Proc.\ CVPR}, pp.~16699--16708, 2024.

\bibitem{sen2008collective} P.~Sen, G.~Namata, M.~Bilgic, L.~Getoor, B.~Gallagher, and T.~Eliassi-Rad, ``Collective classification in network data,'' \textit{AI Mag.}, vol.~29, no.~3, pp.~93--106, 2008.

\end{thebibliography}

\end{document}
